\section{2D Benchmark Results}
\label{sec:results_2d}

This section presents the fixed-$k$ sweep results for the three canonical 2D
domains.
The rectangle, circle, and L-shape each introduce a different geometric
challenge.
The rectangle has uniform boundary normals and no gradient concentration.
The circle has curved boundary normals.
The L-shape has a reentrant corner that concentrates the temperature gradient.
All three NAS-selected architectures were trained at every skip factor from
$k=1$ to $k=5$, giving 45 configurations in total.
Each configuration was run with ten independent random seeds.

Table~\ref{tab:2d_params} lists the physical and training parameters used
across all 2D benchmark runs.

\begin{table}[H]
\centering
\caption{Physical and training parameters for the 2D benchmark simulations.}
\label{tab:2d_params}
\small
\begin{tabular}{llc}
\toprule
Parameter & Symbol & Value \\
\midrule
Thermal conductivity        & $k_T$               & $150\,\mathrm{W/(mK)}$ \\
Volumetric heat capacity    & $\rho c_p$          & $2.4\times10^6\,\mathrm{J/(m^3K)}$ \\
Convection coefficient      & $h$                 & $5000\,\mathrm{W/(m^2K)}$ \\
Initial temperature         & $T_0$               & $540\,^\circ\mathrm{C}$ \\
Water (quench) temperature  & $T_w$               & $20\,^\circ\mathrm{C}$ \\
\midrule
Interior collocation points & $n_\mathrm{domain}$ & 1\,500 \\
Boundary collocation points & $n_\mathrm{bc}$     & 300 \\
Adam epochs (Bayesian/TPE)  & ---                 & 800 \\
Adam epochs (NSGA-II/III)   & ---                 & 1\,500 \\
L-BFGS iterations per window& ---                 & 50 \\
\bottomrule
\end{tabular}
\end{table}

Table~\ref{tab:summary_2d_mae} gives the ten-seed mean~$\pm$~std for all three
domains and architectures.
Full per-seed tables are in Appendix~\ref{app:full_results}.

\begin{table}[H]
\centering\small\setlength{\tabcolsep}{4pt}
\caption{Ten-seed mean-window MAE ($^\circ\mathrm{C}$) summary for all 2D canonical domains.
\cellcolor{RoyalBlue!20}\phantom{x} within $10\,^\circ\mathrm{C}$;
{\color{gray!60!black}grey} = above threshold.}
\label{tab:summary_2d_mae}
\begin{tabular}{@{}llccccc@{}}
\toprule
Domain & Architecture
  & \makecell{$k{=}1$\\(0\%)} & \makecell{$k{=}2$\\(50\%)} & \makecell{$k{=}3$\\(65\%)} & \makecell{$k{=}4$\\(75\%)} & \makecell{$k{=}5$\\(80\%)} \\
\midrule
\multirow{3}{*}{Rectangle 2D}
  & Bayesian/TPE & \cellcolor{RoyalBlue!20}$2.12\pm0.23$ & \cellcolor{RoyalBlue!20}$2.81\pm0.30$ & \cellcolor{RoyalBlue!20}$3.19\pm0.57$ & \cellcolor{RoyalBlue!20}$3.69\pm0.84$ & \cellcolor{RoyalBlue!20}$4.03\pm0.66$ \\
  & NSGA-II      & \cellcolor{RoyalBlue!20}$3.44\pm0.49$ & \cellcolor{RoyalBlue!20}$5.40\pm0.82$ & \cellcolor{RoyalBlue!20}$6.76\pm1.14$ & \cellcolor{RoyalBlue!20}$8.33\pm1.65$ & \cellcolor{RoyalBlue!20}$9.66\pm1.30$ \\
  & NSGA-III     & \cellcolor{RoyalBlue!20}$3.52\pm0.48$ & \cellcolor{RoyalBlue!20}$5.11\pm0.70$ & \cellcolor{RoyalBlue!20}$6.10\pm0.45$ & \cellcolor{RoyalBlue!20}$7.95\pm0.89$ & \cellcolor{RoyalBlue!20}$9.28\pm1.64$ \\
\midrule
\multirow{3}{*}{Circle 2D}
  & Bayesian/TPE & \cellcolor{RoyalBlue!20}$2.75\pm0.33$ & \cellcolor{RoyalBlue!20}$3.58\pm1.17$ & \cellcolor{RoyalBlue!20}$4.93\pm1.95$ & \cellcolor{RoyalBlue!20}$6.13\pm3.36$ & {\color{gray!60!black}$10.01\pm5.16$} \\
  & NSGA-II      & \cellcolor{RoyalBlue!20}$4.89\pm1.50$ & \cellcolor{RoyalBlue!20}$7.58\pm2.86$ & \cellcolor{RoyalBlue!20}$7.34\pm3.63$ & \cellcolor{RoyalBlue!20}$9.31\pm4.61$ & {\color{gray!60!black}$14.97\pm7.81$} \\
  & NSGA-III     & \cellcolor{RoyalBlue!20}$4.01\pm1.16$ & \cellcolor{RoyalBlue!20}$5.83\pm2.41$ & \cellcolor{RoyalBlue!20}$6.82\pm3.55$ & \cellcolor{RoyalBlue!20}$9.45\pm5.96$ & {\color{gray!60!black}$11.33\pm7.58$} \\
\midrule
\multirow{3}{*}{L-Shape 2D}
  & Bayesian/TPE & \cellcolor{RoyalBlue!20}$2.29\pm0.24$ & \cellcolor{RoyalBlue!20}$3.41\pm0.37$ & \cellcolor{RoyalBlue!20}$3.69\pm0.66$ & \cellcolor{RoyalBlue!20}$4.74\pm1.01$ & \cellcolor{RoyalBlue!20}$5.41\pm1.06$ \\
  & NSGA-II      & \cellcolor{RoyalBlue!20}$3.50\pm0.24$ & \cellcolor{RoyalBlue!20}$6.61\pm0.89$ & \cellcolor{RoyalBlue!20}$8.53\pm0.99$ & {\color{gray!60!black}$11.33\pm1.98$} & {\color{gray!60!black}$14.37\pm2.10$} \\
  & NSGA-III     & \cellcolor{RoyalBlue!20}$2.87\pm0.46$ & \cellcolor{RoyalBlue!20}$5.13\pm0.65$ & \cellcolor{RoyalBlue!20}$7.30\pm1.47$ & \cellcolor{RoyalBlue!20}$9.88\pm1.90$ & {\color{gray!60!black}$11.29\pm3.04$} \\
\bottomrule
\end{tabular}
\end{table}

All architectures reproduce the global cooling pattern correctly at their
best configuration.
The hardest regions are near the boundary and at the L-shape corner, where
gradients are steepest.
Temperature field comparisons for all architectures and domains are in
Appendix~\ref{app:temp_fields} (Figure~\ref{fig:app_2d_bestk}).

Prediction error was highest in the early windows across all cases. This is
consistent with the steep temperature gradient at the start of quenching. As
the simulation progressed and the temperature field approached uniformity,
prediction difficulty decreased and the per-window error fell. Bayesian/TPE
showed the most stable window-to-window behaviour across all three 2D domains.
NSGA-II and NSGA-III had slightly higher variance, particularly on the L-shape
domain where the reentrant corner concentrates the temperature gradient.

Figure~\ref{fig:2d_seed_bars} summarises the mean-window MAE across ten
independent seeds for all 2D benchmark configurations. Error bars show one
standard deviation. Full numerical values are provided in the consolidated
MAE tables in Appendix~\ref{app:full_results}
(Tables~\ref{tab:cons_rectangle_2d}--\ref{tab:cons_lshape_2d}).

\begin{figure}[H]
\centering
\includegraphics[width=\textwidth]{seed_bar_2d.png}
\caption{Seed-to-seed variability in mean-window MAE for all 2D domains.
Line: ten-seed mean; band: $\pm$one standard deviation.}
\label{fig:2d_seed_bars}
\end{figure}
\begin{figure}[H]
\centering
\includegraphics[width=\textwidth]{l2_all_k_2d.png}
\caption{Per-window $L_2$ relative error over the 30-second quench for all 2D domains.
Line: ten-seed mean; band: $\pm$one standard deviation. Colour encodes $k$.}
\label{fig:2d_l2_curve}
\end{figure}
